\chapter*{Preface}
\addcontentsline{toc}{chapter}{Preface}

This book is organised around one sentence: \emph{a shape is kinematics; putting in an action produces dynamics.}
Every chapter first shows a shape, then states the action, then derives the equations, and finally points to a theorem that has been checked by the Lean~4 proof assistant.

The shape that runs through the book is the \emph{dumbbell}: two round caps joined by a cylindrical neck.
With the Einstein--Yang--Mills action the caps become closed radiation universes and the neck becomes the throat of an extremal black hole, $\mathrm{AdS}_2\times S^2$.
Both have the same size $\ell$, fixed by the gauge coupling through $\ell^2=4\pi G/g^2$.
Before we get there we build everything we need from quaternions: rotations, spin, the three-sphere, Lorentz transformations, and the geometry of relativistic momentum.

The intended reader is an undergraduate in the second to fourth year who knows linear algebra, multivariable calculus and the first courses on mechanics and electromagnetism.
General relativity, quantum field theory and group theory are not assumed; we build the parts we use.

\section*{How to read the status tags}
Every result carries one of the following tags, taken directly from the claim ledger of the Lean repository.

\begin{center}
\begin{tabular}{@{}ll@{}}
\proved & checked in Lean from its stated hypotheses \\
\external & uses a theorem from the literature, cited in the text \\
\numerical & a numerical computation, not checked by Lean \\
\negative & an idea that was tried and does not work \\
\open & not solved here (and in most cases not solved anywhere)
\end{tabular}
\end{center}

Physical readings (``this is a gravitational wave'') are kept separate from theorems and are called interpretations.
Exercises stay within what Lean has checked; each exercise names the theorem that answers it.

\chapter*{A guide for the specialist}
\addcontentsline{toc}{chapter}{A guide for the specialist}

This section is written for readers who already know general relativity and gauge theory and want to know, in two pages, what is in this book that they have not seen before, and how far it can be trusted.

\section*{What has been checked}
The book is a companion to the Lean~4 repository \emph{Cyclic Fueter Theory} (Zenodo concept DOI \href{https://doi.org/10.5281/zenodo.23448911}{10.5281/zenodo.23448911}, version 0.6.0).
At that version the repository contains 1873 theorem declarations in 171 modules.
An axiom audit of 2421 entries (all public theorems and the index aliases) found no \texttt{sorry}; 2396 entries depend only on Lean's standard axioms \texttt{propext}, \texttt{Classical.choice} and \texttt{Quot.sound}, and the rest on a subset of them.
Lean verifies each theorem from its declared hypotheses.
It does not verify that the hypotheses describe nature, and results quoted from the literature are tagged \external.

\section*{Claims of novelty}
The following statements are, to the author's knowledge, not stated in this form in the literature the author has consulted.
The claim is deliberately weak: no systematic review has been carried out, no specialist has confirmed it, and several items are special cases or rearrangements of known results.
They are listed so that they can be checked.

\begin{enumerate}[leftmargin=*,label=(\arabic*)]
\item \textbf{Linear rigidity} (Chapter~19). Every linear perturbation of the Einstein--Yang--Mills Tolman universe that preserves spherical symmetry, the magnetic ansatz and conformal flatness is generated by a conformal Killing vector of the Einstein static universe, so it is pure gauge.
\item \textbf{Shape of the dumbbell} (Chapter~16). Time-symmetric, conformally flat slices satisfy the Bogomolny-type inequality $r'^2\le W^2$. The case $W=\pm r'$ saturates it and gives $r^2+\ell^2r'^2=\ell^2$. The cap is the regular solution of this equation and the Bertotti--Robinson neck is its singular solution, the envelope of the caps. The junction is fixed by continuity alone, is $C^1$ and never $C^2$, and the Hamiltonian constraint admits exactly the two values $0$ and $-1/\ell$ of $r''$ there.
\item \textbf{Null junction} (Chapters~16--17). A timelike junction between cap and neck is impossible, and the jump $[W_u]=1/(\sqrt2\,\ell)$ is transported unchanged along the null fronts.
\item \textbf{Weyl invariant} (Chapter~16). In this sector the Weyl invariant equals $3\bigl(2m/r-\ell^2(1-W^2)^2/r^2\bigr)$, so conformal flatness says that the Misner--Sharp mass equals the magnetic self-energy on every sphere. This is a special case of the known relation between Weyl curvature and the Misner--Sharp mass.
\item \textbf{One pair of constraints} (Chapter~23). The same Hamiltonian and momentum constraints reduce to the Friedmann constraint on the cap and to a flat two-dimensional Klein--Gordon equation on the Kantowski--Sachs side, where the neck $r=\ell$ is a fold of minisuperspace, invariant under $r\mapsto\ell^2/r$. The neck length is canonically conjugate to $-r\dot r$. Minisuperspace quantisation of Reissner--Nordstr\"om interiors is known; the reading in terms of the dumbbell is what is claimed.
\item \textbf{Constant-area data} (Chapter~26). Time-symmetric data of constant areal radius $r_0$ obey $W'^2=(1-\ell^2(1-W^2)^2/r_0^2)/(2\ell^2)$, with the lump $W=\sqrt2/\cosh(x/\ell)$ at $r_0=\ell$. Only the neck keeps its area at the next instant.
\item \textbf{Weak gravity conjecture} (Chapter~26). With an adjoint Higgs field the electric condition for the $W$ boson and the magnetic condition for the BPS monopole coincide, $4\pi Gv^2\le1$. Together with the stability bound of the neck they give $3/4\le(m_W\ell)^2\le1$, and the neck exists only for $\Lambda\le g^2/(16\pi G)$.
\item \textbf{Formalisation} (whole book). The spherically symmetric Einstein--Yang--Mills--Dirac sector is machine-checked from one reduced action through the constraints, the known limits and the thermodynamic relations.
\end{enumerate}

\section*{What is not claimed}
The cap and neck solutions themselves, the reduced action, the canonical form and the mass function, JT gravity near the neck, black-hole thermodynamics, the sphaleron zero mode, and the weak gravity and Festina Lente conjectures are standard.
The general solution of the Wheeler--DeWitt equation with matter, the quantisation of gravitational waves beyond the linear level, and any derivation of the fine-structure constant are \emph{not} obtained; Chapter~27 lists them as open problems.

\section*{Acknowledgements}
The Lean proofs, numerical checks and parts of the text were prepared with the assistance of a large language model (Claude, Anthropic); the author has checked the results and takes responsibility for the content.
